% TOP_PROBLEM: {"id":30007530,"problem_number":"LOCAL-30007530","title":"Monetary policy and inequality","category_id":21,"category":{"id":21,"name":"economics","display_name":"Economics","description":"Open economic theory statements, published research agendas, and proposed scoped specifications; question provenance and historical priority are qualified per record.","slug":"economics","order_index":21,"created_at":"2026-10-08T00:00:00Z"},"status":"research_question","external_url":null,"legacy_ids":[],"published":false,"statement":"Characterize welfare-maximizing finite-horizon monetary policy with fixed fiscal financing, incomplete insurance, and explicit household welfare weights.","economics":{"original_id":19,"field":"Inflation and monetary policy","kind":"design","formalization_basis":"adapted_research_question","source_status":"explicit_open","priority_status":"earliest_found","priority_note":"This is the earliest explicit relevant statement verified in this audit, not a claim of original priority; older literature and earlier versions may contain antecedents.","earliest_verified_reference":{"authors":["Adrien Auclert","Matthew Rognlie","Ludwig Straub"],"title":"Fiscal and Monetary Policy with Heterogeneous Agents","year":2025,"url":"https://web.stanford.edu/~aauclert/annual_review_hank.pdf","doi":"10.1146/annurev-economics-091624-044646","locator":"Section 5, pp. 19–21 of the April 2025 author manuscript; “Nonlinearities” and “Optimal policy”","evidence_summary":"The review says empirical size and sign asymmetries remain debated. Its optimal-policy discussion identifies unresolved computation and equity–efficiency issues and warns that unrestricted monetary–fiscal Ramsey problems with standard preferences need not have a well-defined steady state."},"current_reference":{"authors":["Adrien Auclert","Matthew Rognlie","Ludwig Straub"],"title":"Fiscal and Monetary Policy with Heterogeneous Agents","year":2025,"url":"https://web.stanford.edu/~aauclert/annual_review_hank.pdf","doi":"10.1146/annurev-economics-091624-044646","locator":"Section 5, pp. 19–21 of the April 2025 author manuscript; “Nonlinearities” and “Optimal policy”","evidence_summary":"The review says empirical size and sign asymmetries remain debated. Its optimal-policy discussion identifies unresolved computation and equity–efficiency issues and warns that unrestricted monetary–fiscal Ramsey problems with standard preferences need not have a well-defined steady state."},"formalization_note":"The review explicitly identifies optimal heterogeneous-agent policy as unfinished and warns about unrestricted steady-state problems. This restricted finite-horizon monetary target avoids that pathology."},"statusQualification":"Published question or research agenda verified at the cited locator. The mathematical specification is adapted for this catalogue and includes additional assumptions; source publication does not establish first-ever priority or certify this exact model remains unsolved. The review explicitly identifies optimal heterogeneous-agent policy as unfinished and warns about unrestricted steady-state problems. This restricted finite-horizon monetary target avoids that pathology.","statusReviewedAt":"2026-10-08"}
% FIRST_POSTED: 2026-10-08
% ENTRY_KIND: statement_only
% !TeX program = lualatex
\documentclass[11pt]{article}
\usepackage{fontspec}
\usepackage[a4paper,margin=25mm]{geometry}
\usepackage{amsmath,amssymb,mathtools}
\usepackage{xurl}
\usepackage[hidelinks]{hyperref}
\newcommand{\sref}[2]{\hyperlink{source:#1:#2}{[S#2]}}
\setlength{\emergencystretch}{3em}
\begin{document}
% problemId:30007530
\section*{Monetary policy and inequality}\label{30007530}
\subsection{Definitions and mathematical statement}
Fix a finite-horizon heterogeneous-agent monetary economy with incomplete insurance, idiosyncratic earnings risk, nominal assets, sticky prices, and horizon \(T\). Household types \(i\) have utility \(u_i(c,\ell)\), social weights \(\omega_i\ge0\), and discount factor \(\beta\). Initial distributions and a terminal resource and debt rule are given. Government transfers and taxes follow a fixed feasible rule, so that the monetary authority controls only a policy-rate path \(a=(i_0,\ldots,i_T)\) subject to an effective lower bound and an information constraint.
For every feasible path inducing an equilibrium, define
\[
W(a)=E\sum_{t=0}^T\beta^t\sum_i\omega_i u_i(c_{it}(a),\ell_{it}(a)).
\]
Characterize a maximizing path and the marginal welfare terms associated with nominal redistribution, unemployment risk, price dispersion, and constrained consumption. Determine whether an implementable rule based on inflation, employment, and measured distributional states approximates this path within a declared welfare tolerance. Identify the exposure and preference parameters needed for those comparisons, and report sensitivity to social weights rather than treating them as estimated facts. The finite horizon and fixed fiscal rule avoid assuming the existence of an unrestricted infinite-horizon Ramsey steady state. The research target is a quantitatively disciplined equity–efficiency policy assessment, not a claim that reducing any inequality statistic always improves monetary policy.

\par\textbf{Formulation and scope.} The review explicitly identifies optimal heterogeneous-agent policy as unfinished and warns about unrestricted steady-state problems. This restricted finite-horizon monetary target avoids that pathology.

\par\textbf{Open-question provenance.} An explicit question or research agenda was verified at the source locator. This does not by itself certify that every later version remains unresolved \sref{30007530}{1}.

\par\textbf{Historical priority.} This is the earliest explicit relevant statement verified in this audit, not a claim of original priority; older literature and earlier versions may contain antecedents.

\subsection{Short English statement}
Characterize welfare-maximizing finite-horizon monetary policy with fixed fiscal financing, incomplete insurance, and explicit household welfare weights.

\subsection{Sources}
\begin{itemize}\raggedright
\item[\textnormal{[S1]}]\hypertarget{source:30007530:1}{} Adrien Auclert, Matthew Rognlie, Ludwig Straub. 2025. Fiscal and Monetary Policy with Heterogeneous Agents. Section 5, pp. 19–21 of the April 2025 author manuscript; “Nonlinearities” and “Optimal policy”.
\url{https://web.stanford.edu/~aauclert/annual_review_hank.pdf}.
DOI: 10.1146/annurev-economics-091624-044646.
{\small\textit{Earliest verified question/agenda reference; historical priority qualified below. The review says empirical size and sign asymmetries remain debated. Its optimal-policy discussion identifies unresolved computation and equity–efficiency issues and warns that unrestricted monetary–fiscal Ramsey problems with standard preferences need not have a well-defined steady state.}}
\end{itemize}
\end{document}
